%======================================================================
\chapter{Der Tensor Power Flow}
\label{ch:tpf}
%======================================================================

In \cref{ch:grundlagen} wurden die elektrischen und methodischen Grundlagen des NR-Verfahrens und der Fixpunkt-Iteration eingeführt. 
Fortlaufend wird auf dieser Basis die Tensorisierung der FPI eingeführt, welches den Tensor Power Flow (TFP) bildet.
Diese Mehrdimensionale Formulierung orientiert sich nach Salazar Duque et al. \cite{Salazar2024}.
Dabei wird der strukturelle Vorteil des FPI, einer konstanten Systemmatrix, hier ausgenutzt.
Daraus resultiert, dass das Verfahren Tausende von Lastflüssen gleichzeitig in jeder Iteration berechnen kann und somit für reine PQ-Netze eine erhebliche Beschleunigung erzielt werden kann.\\

Zunächst wird der Algorithmus (\cref{sec:tpf_algorithmus}) eingeführt, welcher darauf folgend in seiner Dense-Formulierungen vertieft wird. 
Die Sparse-Formulierung wird für diese Ausarbeitung nicht verwendet, da sie erst ab Netzgrößen von über 1000 Knoten effizienter sein kann.
Der Fokus bleibt auf der Analyse der einer äußeren Schleife.
Das Kapitel wird damit abgeschlossen, dass der Zentrale Nachteil dieser Berechnungsvorschrift erklärt wird: Die fehlende Möglichkeit PV-Knoten direkt miteinzubeziehen. 
Damit wird die in Kapitel 4 vorgestellte Methode zur Berechnung von PV-Knoten mit dem TPF motiviert.

%======================================================================
\section{Grundalgorithmus}
\label{sec:tpf_algorithmus}
%======================================================================

Das in \cref{sec:fpi} zugrunde gelegte FPI-Verfahren dient als Ausgangspunkt für den TPF. 
Es wird weiterhin die Nomenklatur der eingeführten Netzbeschreibung \eqref{eq:admittance_block} übernommen.
Dabei wird zusätzlich zu den $b$ Lastknoten $\phi$ für die Anzahl der Phase verwendet.
Das Multiplikat der beiden bildet die Anzahl der Busphasen $b\phi$. 
Die zusammengefasste Iterationsvorschrift auf Basis von \eqref{eq:fpi_iteration} ist durch


\begin{equation}
	\mathbf{v}_{(n+1)} = \mathbf{F} \cdot \mathbf{v}_{(n)}^{*\eleminv} + \mathbf{w}
	\label{eq:tpf_iteration}
\end{equation}

gegeben, wobei $\mathbf{v}_{(n)}^{*\eleminv} \in \C^{b\phi \times 1}$ die komponentenweise, inverse, komplex konjugierte Spannung 


\begin{equation}
	\left(\mathbf{v}^{*\eleminv}\right)_i = \frac{1}{v_i^*} \quad  i = 1, \ldots, b\phi
	\label{eq:eleminv_def}
\end{equation}

darstellt.
Dabei ist $n$ die aktuelle Iteration der Spannungsberechnung und die Hilfsmatrizen werden ähnlich wie in \cref{sec:fpi} formuliert mit


\begin{align}
	\mathbf{A} &= \diag(\boldsymbol{\alpha_P} \had \mathbf{s}^*)  &\in \C^{b\phi \times b\phi}\label{eq:tpf_A} \\
	\mathbf{B} &= \diag(\boldsymbol{\alpha_Z} \had \mathbf{s}^*) + \mathbf{\Ydd} &\in \C^{b\phi \times b\phi} \label{eq:tpf_B}\\
	\mathbf{c} &= \Yds \mathbf{v}_s + \boldsymbol{\alpha_I} \had \mathbf{s}^* &\in \C^{b\phi \times 1} \label{eq:tpf_c}\\
	\mathbf{F} &= -\mathbf{B}^{-1}\mathbf{A} &\in \C^{b\phi \times b\phi} \label{eq:tpf_F}\\
	\mathbf{w} &= -\mathbf{B}^{-1}\mathbf{c} &\in \C^{b\phi \times 1}\label{eq:tpf_w}.
\end{align}

Die Vektoren und Matrizen $\mathbf{A}$, $\mathbf{B}$ und $\mathbf{c}$ sind bei gegebenen Slack-Spannungen und Leistungen Konstant über alle Iterationen. 
Daraus resultiert, dass die Systemmatrix $\mathbf{F}$ und der Offset $\mathbf{w}$ nur ein einziges mal aufgestellt werden müssen.\\


Eine kompaktere Darstellung der Iterationsvorschrift lässt sich in dem gängigen Analysefall der konstanten Last bei $\alpha_P = 1$, $\alpha_I = \alpha_Z = 0$ formulieren~\cite{Kundur1994,Kersting2017,IEEETF1995,VanCutsem1998,Bolognani2016}, wobei die vereinfachten Hilfsgrößen

\begin{align}
	\mathbf{B} &= {\Ydd} \label{eq:tpf_B_konstP}\\
	\mathbf{F} &= -{\Ydd}^{-1} \diag(\mathbf{s}^*) = -{\ZB} \diag(\mathbf{s}^*) \label{eq:tpf_F_konstP}\\
	\mathbf{w} &= -{\ZB} \Yds \mathbf{v}_s \label{eq:tpf_w_konstP}
\end{align}


mit $\mathbf{\ZB} = \mathbf{\Ydd}^{-1}$ als Bus-Impedanzmatrix den vereinfachen.
Der TPF vereinfacht sich entsprechend zu 
\begin{equation}
	\mathbf{v}_{(n+1)} = \ZB \cdot \left(\vect{s}^* \had \mathbf{v}_{(n)}^{*\eleminv}\right) + \mathbf{w}
	\label{eq:tpf_konstP}
\end{equation}

wobei die Abhängigkeit zu der Knotenleistung durch den Klammerausdruck leichter nachzuvollziehen ist.
Genau dieser entspricht den konjugierten Knotenströmen $\mathbf{i}^* = \mathbf{s}^* \oslash \mathbf{v}_{(n)}^*$.
Zusammen mit der Bus-Impedanzmatrix entspricht der Term $\ZB \cdot \vect{i}^*$ dem Spannungsabfall der Knoten, während $\vect{w}$ den Zusatz der Slack-Spannungen beschreibt.
Zur besseren Übersichtlichkeit wird fortan diese Darstellung verwendet.
In der Implementierung wird jedoch weiterhin der allgemeine ZIP-Fall berücksichtigt.


%======================================================================
\section{Dense-Formulierung}
\label{sec:tpf_formulierungen}

Zeitreihensimulationen, probabilistischen Lastflüsse, sowie das maschinelle Lernen sind Anwendungen, wo eine Vielzahl an Lastflüssen berechnet werden müssen.
Dadurch bedingt liegen solche Daten oft in Form von mehrdimensionalen Datenstrukturen vor \cite{Salazar2024, Sandoval.2020, SandovalGuzman.2022, Zhang.2015}.
Die Lastdimensionen können beispielsweise die Anzahl der Zeitschritte $t$, der Lastszenarien $r$, der zu durchführenden Experimente $p$ und auch die Anzahl der Busphasen $b\phi$ sein. 
Diese Leistungsdaten können kompakt in der Form eines Tensors 

\begin{equation}
	\tens{S} \in \C^{\ldots \times p \times r \times t \times b\phi}
	\label{eq:tensor_S}
\end{equation}

aufgelistet werden. 
Die Tensordimension setzt sich aus den verschiedenen Anwendungsarten mit deren Anzahl zu $\ldots \times p \times r \times t \times b\phi$ zusammen.
bis auf die Busphasen können die nichtphysikalischen Dimensionen zu einer zusammengelegt werden (\emph{reshape}).
Dadurch wird die Struktur auf $\tau \times b\phi$ vereinfacht werden wobei $\tau = \ldots \cdot p \cdot r \cdot t$ die restlichen Dimensionen zusammenhängt~\cite{Salazar2024}.
Diese \emph{reshape}-Operation 

\begin{equation}
	\dot{\mathbf{S}} \in \C^{b\phi \times \tau}
	\label{eq:reshape}
\end{equation}

generiert den Leistungstensor $\dot{\mathbf{S}}$, welcher in jeder Spalte $i \in \left[0, \tau \right]$ die Leistung eines Szenarios beinhaltet und die Zeile, wie zuvor den Bus (ggf. die Phase) des Leistungsflusses kennzeichnet.
Die restlichen Matrizen folgen der selben Konvention.
Wird der FPI aus \eqref{eq:tpf_iteration} tensorisiert dargestellt, folgt zunächst

\begin{equation}
	\tens{V}_{(n+1)} = \tens{F} \cdot \tens{V}_{(n)}^{*\eleminv} + \tens{W}
	\label{eq:tensor_fpi}
\end{equation}
wobei $\tens{V} \in \C^{\ldots \times p \times r \times b\phi \times t}, \tens{F} \in \C^{\ldots \times p \times r \times t \times b\phi \times b\phi}, \tens{W} \in \C^{\ldots \times p \times r \times b\phi \times t}$ die korrespondierenden Tensorgrößen von $\mathbf{V}, \mathbf{F}$ und $\mathbf{w}$ sind.
Unter Verwendung des konstanten Lastfalls und einer konstanten Netztopologie über alle Szenarien und Iterationen ist der Tensor $\tens{F}$ eine reine Wiederholung der Bus-Impedanzmatrix $\ZB$. 
Ferner ist auch der Offset-Tensor $\tens{W}$ ein Reihe an Wiederholungen des Slack-Einflusses, welcher für diese Ausarbeitung ebenso über alle Szenarien konstant bleiben soll und somit die Form $\dot{W} = [\mathbf{w}, \mathbf{w}, \ldots, \mathbf{w}] \in \C^{b\phi \times \tau}$ hat.
Die Iteration erfolgt nach dem \emph{reshape} dem Ausdruck


\begin{equation}
	\dotV_{(n+1)} = \ZB \cdot \underbrace{\left(\dotS^* \had \dotV_{(n)}^{*\eleminv}\right)}_{\in \C^{b\phi \times \tau}} + \dot{W},
	\label{eq:dense_matmul}
\end{equation}

welcher die selbe Struktur hat wie \eqref{eq:tpf_konstP} \cite{Salazar2024}.
Dabei kann in diesem Fall die selbe Bus-Impedanzmatrix verwendet werden.
Vor der Tensorformlierung liegt eine Vektor-Matrix Multiplikation zwischen der Impedanzmatrix $\ZB$ und dem Stromvektor $\vect{s}^* \had \mathbf{v}_{(n)}^{*\eleminv}$ vor. 
Nach der Tensorisierung und dem \emph{reshape} wird diese Operation zu einem Matrix-Matrix-Produkt zwischen $\ZB$ und $\dotS^* \had \dotV_{(n)}^{*\eleminv}$.
Hierbei werden alle Leistungsvektoren der Lastflüsse berücksichtigt.
Der nachfolgende Algorithmus zeigt Dense-Formulierung nach Salazar Duque \cite{Salazar2024}.
\todo{hier weiter}




\begin{algorithm}[H]
	\caption{Tensor Power Flow. Dense-Formulierung nach \cite{Salazar2024}}
	\label{alg:tpf_dense}
	\begin{algorithmic}[1]
		\REQUIRE $\dotS \in \C^{b\phi \times \tau}$, $\ZB \in \C^{b\phi \times b\phi}$, $\vect{w} \in \C^{b\phi \times 1}$, Toleranz $\varepsilon$, max.\ Iterationen $K$
		\ENSURE $\dotV_n \in \C^{b\phi \times \tau}$, Iterationszahl $n$
		\STATE $\dotV_0 = \mathbf{1} + \jj\mathbf{0} \in \C^{b\phi \times \tau}$; \quad $n = 0$; \quad $\text{tol} = \infty$
		\WHILE{$\text{tol} \geq \varepsilon$ \textbf{und} $n < K$}
		\FOR{$i = 1$ \textbf{bis} $\tau$}
		\STATE $\dotV_{(n+1)}[i] = \ZB \left(\dotS[i] \had \dotV_{(n)}^{\eleminv}[i]\right)^* + \vect{w}$
		\ENDFOR
		\STATE $\text{tol} = \max\left(\|s_{(n+1)} - s_{(n)}\|_2\right)$; \quad $n = n + 1$
		\ENDWHILE
		\RETURN $\dotV_{(n)}$, $n$
	\end{algorithmic}
\end{algorithm}



%\subsection{Sparse-Formulierung}
%\label{subsec:tpf_sparse}
%
%Ein Nachteil der dense Formulierung ist, dass die Matrix $\ZB = \Ydd^{-1}$ vollbesetzt (dense) ist, obwohl die Admittanzmatrix $\Ydd$ dünnbesetzt (sparse) ist. Für große Netze ($b\phi > 1000$) kann die Speicherung von $\ZB$ problematisch werden. Salazar Duque et al.\ \cite{Salazar2024} schlagen daher eine sparse Formulierung vor, die die Inversion von $\Ydd$ vermeidet.
%Die Iteration \eqref{eq:tpf_iteration} wird umformuliert zu:
%\begin{equation}
%	\tens{M} \cdot \tens{V}_{(n+1)} = \tens{V}_{(n)}^{*\eleminv} + \tens{H}
%	\label{eq:sparse_tensor}
%\end{equation}
%
%wobei:
%\begin{align}
%	\tens{M} &= -\vect{A}^{\eleminv} \vect{B} \label{eq:sparse_M}\\
%	\tens{H} &= \vect{A}^{\eleminv} \vect{C} \label{eq:sparse_H}
%\end{align}
%
%mit $\vect{A}$, $\vect{B}$ und $\vect{C}$ gemäß \eqref{eq:tpf_A}--\eqref{eq:tpf_c}. Da $\vect{A}$ eine Diagonalmatrix ist, überträgt sich die Dünnbesetztheit von $\Ydd$ auf $\tens{M}$.
%
%Die Reshaped-Form von \eqref{eq:sparse_tensor} ergibt ein lineares Gleichungssystem der Struktur $\vect{A}\vect{x} = \vect{b}$:
%\begin{equation}
%	\underbrace{\dotM}_{\vect{A}} \cdot \underbrace{\dotV_{(n+1)}}_{\vect{x}} = \underbrace{\dotV_{(n)}^{*\eleminv} + \dot{H}}_{\vect{b}}
%	\label{eq:sparse_system}
%\end{equation}
%
%Die Matrix $\dotM \in \C^{(\tau \cdot b\phi) \times (\tau \cdot b\phi)}$ ist die Block-Diagonalverkettung der Submatrizen:
%\begin{equation}
%	\dotM = \begin{bmatrix}
%		\vect{M}_1 & 0 & \cdots & 0 \\
%		0 & \vect{M}_2 & \cdots & 0 \\
%		\vdots & & \ddots & \vdots \\
%		0 & 0 & \cdots & \vect{M}_\tau
%	\end{bmatrix}
%	\label{eq:block_diagonal}
%\end{equation}
%
%Die Vektoren $\dotV_{(n+1)}$, $\dotV_{(n)}^{*\eleminv}$ und $\dot{H}$ werden vertikal angeordnet. Das System \eqref{eq:sparse_system} wird iterativ mittels eines sparse Direktlösers berechnet.

%\begin{algorithm}[H]
%	\caption{Tensor Power Flow -- Sparse-Formulierung nach \cite{Salazar2024}}
%	\label{alg:tpf_sparse}
%	\begin{algorithmic}[1]
%		\REQUIRE $\dotS \in \C^{b\phi \times \tau}$, $\Ydd \in \C^{b\phi \times b\phi}$ (sparse), $\Yds \in \C^{b\phi \times \phi}$ (sparse), Toleranz $\varepsilon$, max.\ Iterationen $K$
%		\ENSURE $\dotV_n \in \C^{b\phi \times \tau}$, Iterationszahl $n$
%		\STATE Berechne $\dotM$ und $\dot{H}$ (Block-Diagonal-Aufbau gemäß \eqref{eq:block_diagonal})
%		\STATE $\dotV_0 = \mathbf{1} + \jj\mathbf{0} \in \C^{(\tau \cdot b\phi) \times 1}$; \quad $n = 0$; \quad $\text{tol} = \infty$
%		\WHILE{$\text{tol} \geq \varepsilon$ \textbf{und} $n < K$}
%		\STATE $\dotV_{(n+1)} = \text{solve}\left(\dotM, \; \dotV_{(n)}^{*\eleminv} + \dot{H}\right)$ \COMMENT{Sparse-Löser}
%		\STATE $\text{tol} = \max\left(\|s_{(n+1)} - s_{(n)}\|_2\right)$; \quad $n = n + 1$
%		\ENDWHILE
%		\STATE $\vect{V}_n = \text{reshape}(\dotV_n, \; b\phi \times \tau)$
%		\RETURN $\vect{V}_n$, $n$
%	\end{algorithmic}
%\end{algorithm}

%\subsection{Entscheidender Vorteil der Sparse-Formulierung}
%\label{subsec:sparse_vorteil}
%
%Sparse Direktlöser (z.\,B.\ PARDISO, UMFPACK) lösen lineare Gleichungssysteme in drei Schritten \cite{Davis2006}:
%\begin{enumerate}
%	\item \textbf{Symbolische Analyse:} Bestimmung der Fill-In-Struktur und Permutierung zur Minimierung des Fill-In.
%	\item \textbf{Numerische Faktorisierung:} LU-Zerlegung (oder Cholesky-Zerlegung) der Matrix.
%	\item \textbf{Vorwärts-/Rückwärtssubstitution:} Lösung des faktorisierten Systems.
%\end{enumerate}
%
%Der entscheidende Vorteil der Formulierung \eqref{eq:sparse_system} liegt darin, dass die Systemmatrix $\dotM$ \textit{über alle Iterationen konstant} bleibt und sich nicht ändert. Daher müssen die rechentechnisch aufwändigen Schritte~(1) und~(2) nur \textit{einmal} -- in der ersten Iteration -- durchgeführt werden. In allen folgenden Iterationen wird nur der vergleichsweise günstige Schritt~(3) ausgeführt. \Cref{tab:sparse_schritte} stellt dies dem NR-Verfahren gegenüber.
%
%\begin{table}[H]
%	\centering
%	\caption{Vergleich der sparse Löser-Schritte pro Iteration: NR vs.\ TPF (Sparse).}
%	\label{tab:sparse_schritte}
%	\begin{tabular}{@{}lcc@{}}
%		\toprule
%		\textbf{Schritt} & \textbf{Newton-Raphson} & \textbf{TPF (Sparse)} \\
%		\midrule
%		(1) Symbolische Analyse & Jede Iteration\textsuperscript{*} & Einmal \\
%		(2) Numerische Faktorisierung & Jede Iteration & Einmal \\
%		(3) Vorwärts-/Rückwärtssubstitution & Jede Iteration & Jede Iteration \\
%		\bottomrule
%	\end{tabular}\\[3pt]
%	{\footnotesize \textsuperscript{*} Die symbolische Analyse kann im NR wiederverwendet werden, wenn sich die Sparsity-Struktur nicht ändert. Die numerische Faktorisierung muss jedoch stets neu berechnet werden, da sich die Jacobi-Matrix in jeder Iteration ändert.}
%\end{table}

%%======================================================================
%\section{Parallelisierung auf CPU und GPU}
%\label{sec:tpf_parallel}
%%======================================================================
%
%\subsection{Warum der TPF parallelisierbar ist}
%\label{subsec:tpf_warum_parallel}
%
%Die Parallelisierbarkeit des TPF ergibt sich aus zwei Eigenschaften:
%
%\begin{enumerate}
%	\item \textbf{Konstante Systemmatrix:} $\ZB$ (dense) bzw.\ $\dotM$ (sparse) ist für alle $\tau$ Lastflüsse und über alle Iterationen identisch. Im Gegensatz dazu besitzt beim NR-Verfahren jeder Lastfluss seine eigene, iterationsabhängige Jacobi-Matrix.
%	
%	\item \textbf{Unabhängigkeit der Lastflüsse:} Die $\tau$ Lastflüsse sind untereinander vollständig entkoppelt
%\end{enumerate}
%
%Diese Eigenschaften ermöglichen es, die Berechnung als Matrix-Matrix-Multiplikation $\ZB \cdot \dot{X}$ mit $\dot{X} \in \C^{b\phi \times \tau}$ zu formulieren (vgl.\ \eqref{eq:dense_matmul}). Solche Operationen werden von optimierten BLAS-Bibliotheken (Basic Linear Algebra Subprograms) wie OpenBLAS, Intel MKL oder cuBLAS hocheffizient ausgeführt.
%
%\subsection{CPU-Parallelisierung}
%\label{subsec:tpf_cpu}
%
%Auf modernen Mehrkern-CPUs profitiert die dense Formulierung \eqref{eq:dense_matmul} von den folgenden Mechanismen:
%
%\begin{itemize}
%	\item \textbf{BLAS Level-3 Operationen:} Matrix-Matrix-Multiplikationen (DGEMM/ZGEMM) nutzen automatisch alle verfügbaren CPU-Kerne und optimieren die Cache-Auslastung.
%	\item \textbf{Vektorisierung:} SIMD-Instruktionen (Single Instruction, Multiple Data) verarbeiten mehrere Datenpunkte gleichzeitig.
%	\item \textbf{Effiziente Speicherzugriffe:} Bei einer Matrix-Matrix-Multiplikation werden die Daten im Cache wiederverwendet (hohe Cache-Hit-Rate), im Gegensatz zu $\tau$ einzelnen Matrix-Vektor-Multiplikationen, bei denen $\ZB$ wiederholt geladen werden müsste.
%\end{itemize}
%
%Die sparse Formulierung profitiert von der einmaligen Faktorisierung: PARDISO \cite{Davis2006} parallelisiert sowohl die LU-Zerlegung als auch die Vorwärts-/Rückwärtssubstitution über mehrere Kerne.
%
%\subsection{GPU-Parallelisierung}
%\label{subsec:tpf_gpu}
%
%Grafikprozessoren (GPUs) besitzen tausende einfacher Rechenkerne, die für massiv parallele Datenverarbeitung optimiert sind. Die dense TPF-Formulierung eignet sich besonders für GPUs, da:
%
%\begin{itemize}
%	\item Die Matrix-Matrix-Multiplikation $\ZB \cdot \dot{X}$ durch GPU-optimierte Bibliotheken (cuBLAS, cuPy, PyTorch) mit hohem Durchsatz berechnet werden kann.
%	\item Die element-weisen Operationen ($\had$, $(\cdot)^{*\eleminv}$) trivial parallelisierbar sind.
%	\item Der Speichertransfer (Host $\to$ GPU) nur einmal erfolgen muss, da $\ZB$ und $\vect{w}$ konstant sind.
%\end{itemize}
%
%Allerdings zeigen die Ergebnisse in \cite{Salazar2024}, dass die GPU-Implementierung bei einer geringen Anzahl von Lastflüssen ($\tau \leq 100$) aufgrund des Overhead für den Datentransfer zwischen Host und GPU langsamer ist als die CPU-Implementierung. Erst bei großen $\tau$ (ab ca.\ $\tau > 500$ je nach Netzgröße) überwiegt der Parallelisierungsvorteil der GPU.
%
%\subsection{Rechenzeitvergleich}
%\label{subsec:tpf_ergebnisse}
%
%Salazar Duque et al.\ \cite{Salazar2024} vergleichen die Rechenzeiten der TPF-Formulierungen mit konventionellen Methoden (NR in Sparse-Formulierung, BFS, SAM) anhand synthetischer Netze mit $b\phi \in [9, 5000]$ und $\tau \in [1, 525600]$. Die Experimente wurden auf einem Laptop mit Intel Core i7-7700HQ CPU und NVIDIA Quadro M1200 GPU durchgeführt.
%
%\Cref{tab:tpf_ergebnisse} fasst die Ergebnisse für eine Jahressimulation bei verschiedenen Zeitauflösungen zusammen.
%
%\begin{table}[H]
%	\centering
%	\caption{Rechenzeiten für eine Jahressimulation bei verschiedenen Zeitauflösungen. Werte in Minuten. Daten aus \cite[Table~I]{Salazar2024}.}
%	\label{tab:tpf_ergebnisse}
%	\begin{tabular}{@{}llrrrr@{}}
%		\toprule
%		\textbf{Algorithmus} & \textbf{$b\phi$} & \textbf{@1\,h} & \textbf{@30\,min} & \textbf{@15\,min} & \textbf{@1\,min} \\
%		& & (8.760 PF) & (17.520 PF) & (35.040 PF) & (525.600 PF) \\
%		\midrule
%		BFS & 100 & 1,98 & 3,96 & 7,91 & 118,71 \\
%		NR (Sparse) & 100 & 1,03 & 2,06 & 4,12 & 61,85 \\
%		Tensor (Sparse) & 100 & 0,17 & 0,34 & 0,68 & 10,14 \\
%		SAM & 100 & 0,02 & 0,04 & 0,09 & 1,29 \\
%		Tensor (GPU) & 100 & 0,02 & 0,04 & 0,09 & 1,30 \\
%		\textbf{Tensor (Dense)} & \textbf{100} & \textbf{0,01} & \textbf{0,01} & \textbf{0,02} & \textbf{0,37} \\
%		\midrule
%		BFS & 5000 & 186,49 & 372,99 & 745,97 & -- \\
%		NR (Sparse) & 5000 & 8,77 & 17,54 & 35,08 & 526,23 \\
%		Tensor (Sparse) & 5000 & 2,43 & 4,86 & 9,71 & 145,71 \\
%		SAM & 5000 & 108,09 & 216,18 & 432,36 & -- \\
%		\textbf{Tensor (GPU)} & \textbf{5000} & \textbf{0,92} & \textbf{1,84} & \textbf{3,69} & \textbf{55,34} \\
%		Tensor (Dense) & 5000 & 2,60 & 5,19 & 10,38 & 155,71 \\
%		\bottomrule
%	\end{tabular}
%\end{table}
%
%Aus \cref{tab:tpf_ergebnisse} ergeben sich die wesentlichen Erkenntnisse:
%
%\begin{itemize}
%	\item Für kleine Netze ($b\phi = 100$) ist \textbf{Tensor (Dense)} der schnellste Algorithmus. Die Jahressimulation bei 1-Minuten-Auflösung (525.600 Lastflüsse) benötigt lediglich 22 Sekunden -- ein \textbf{Speedup-Faktor von 164} gegenüber NR (Sparse) mit über einer Stunde Rechenzeit.
%	\item Für große Netze ($b\phi = 5000$) wird \textbf{Tensor (GPU)} zur schnellsten Methode, da die GPU die große Matrix-Matrix-Multiplikation massiv parallelisiert.
%	\item Die FPI-basierten Methoden (SAM, Tensor Dense/Sparse/GPU) sind den konventionellen Methoden (NR, BFS) bei großen $\tau$ durchgehend überlegen.
%	\item Die sparse Formulierung skaliert besser mit der Netzgröße als die dense Formulierung, da sie die Speicherung von $\ZB$ vermeidet.
%\end{itemize}
%
%Die Wahl zwischen dense und sparse Formulierung hängt somit von der Netzgröße ab: Für $b\phi \lesssim 1000$ ist die dense Formulierung (mit oder ohne GPU) vorteilhaft, für größere Netze die sparse Formulierung oder die GPU-Variante.

%======================================================================
\section{Limitation PV-Knoten}
\label{sec:tpf_limitation}
%======================================================================
\label{subsec:tpf_annahme}




Wie bereits eingeführt wurde, setzt der TPF voraus, dass der Leistungstensor $\dot{\mathbf{S}}$ über alle Iterationen und Lastflüsse konstant und bekannt bleibt.
Ist dies nicht gegeben, muss im allgemeinen Fall der Tensor $\tens{F}$ neu berechnet werden.
Wenn Komponenten des Leistungsvektors unbekannt sind, befinden sich gesuchte Größen in der Berechnungsvorschrift und die FPI kann nicht durchgeführt werden.\\

Werden nun PV-Knoten betrachtet, ist genau Situation vorhanden.
Bei einem solchen Knoten ist die Leistung nur in der Wirkleistungskomponente $P$ bekannt, während die Blindleistung $Q$ eine zu berechnende Größe darstellt.
Dadurch kann die Systemmatrix nicht mehr einmalig im voraus berechnet werden und wiederverwendet werden.\\

In der Nachfolgende Methode wird die Blindleistung $Q$ iterativ bestimmt, wodurch die Konstanz der Matrix $\mathbf{A} = \diag(\alpha_P \had \mathbf{s}^*)$ verloren geht. 
Wenn ein konstanter Lastfall mit $\alpha_P = 1$ angenommen wird, ändert sich die Systemmatrix $\ZB$ zwar nicht mehr über die Iterationen, jedoch ist die Formulierung für den Strom $\dotS^* \had \dotV_{(n)}^{*\eleminv}$ weiterhin von der Leistung abhängig und benötigt ebenfalls eine eigenständige Handhabung der unbekannten PV-Blindleistungen.\\

Bislang wurde diese Einschränkung des TPF nicht auf PV-Knoten erweitert. 
Der reine PQ-Knoten-Fall reicht in vielen Fällen aus.
Kleine PV-Anlagen, Windkraftanlagen, Gewerbelasten, Batterien und  Wechselrichter können alle als PQ Knoten dargestellt werden, wobei PV-Knoten bislang zu PQ-Knoten umgewandelt werden, um diese Beschränkung der Methode zu vermeiden~\cite{Salazar2024}.
Im Gegensatz wird mit zunehmender Anzahl von dezentralen Erzeugeranlagen, wie Synchrongeneratoren und PV-Anlagen, eine Option zur Inklusion von PV-Knoten im TPF relevanter~\cite{Ghiani2015,Wang2019}. 






%Die Effizienz des Tensor Power Flow basiert auf einer fundamentalen Annahme: \textit{Der Leistungsvektor $\vect{s}$ ist für alle $\tau$ Lastflüsse vollständig bekannt und bleibt während der gesamten Iteration konstant.}
%
%Diese Annahme ermöglicht:
%\begin{enumerate}
%	\item Die \textbf{Vorausberechnung} der Systemmatrix $\ZB$ (bzw.\ $\dotM$) und des Konstantvektors $\vect{w}$.
%	\item Die \textbf{Wiederverwendung} dieser Größen über alle Iterationen und alle $\tau$ Lastflüsse.
%	\item Die \textbf{Zusammenfassung} aller $\tau$ Lastflüsse in eine einzige Matrix-Matrix-Multiplikation.
%\end{enumerate}
%
%Ohne diese Konstantheit wäre der Tensor-Reshape \eqref{eq:reshape} nicht durchführbar, da jeder Lastfluss seine eigene, sich ändernde Systemmatrix benötigen würde.
%
%%\subsection{Warum PV-Knoten diese Annahme verletzen}
%%\label{subsec:tpf_pv_problem}
%
%Bei PV-Knoten ist die Blindleistung $Q_k$ unbekannt (\cref{subsec:knotentypen}). Der Leistungsvektor am PV-Knoten $k$ lautet:
%\begin{equation}
%	s_k = P_k^{\spec} + \jj \underbrace{Q_k}_{\text{unbekannt}}
%	\label{eq:s_pv_unbekannt}
%\end{equation}
%
%Da $Q_k$ iterativ bestimmt werden muss, um die Spannungsbedingung $|V_k| = |V_k^{\spec}|$ zu erfüllen, ändert sich $s_k$ zwischen den Iterationen. Dies hat folgende Konsequenzen für die TPF-Formulierung:
%
%\begin{enumerate}
%	\item Die Matrix $\vect{A} = \diag(\alpha_P \had \vect{s}^*)$ aus \eqref{eq:tpf_A} ist \textbf{nicht mehr konstant}, da $\vect{s}$ sich ändert.
%	\item Damit sind auch $\vect{F} = -\vect{B}^{-1}\vect{A}$ und $\vect{w} = -\vect{B}^{-1}\vect{c}$ nicht mehr über die Iterationen konstant.
%	\item Im Fall rein konstanter Leistung ($\alpha_P = 1$): Die Iteration $\vect{v}_{(n+1)} = \ZB(\vect{s}^* \had \vect{v}_{(n)}^{*\eleminv}) + \vect{w}$ ändert sich nicht in der Systemmatrix $\ZB$ selbst, aber der Leistungsvektor $\dotS$ im Produkt $\dotS^* \had \dotV_{(n)}^{*\eleminv}$ ist für PV-Knoten iterationsabhängig.
%	\item Im Tensor-Setting: Verschiedene Szenarien $i$ und $j$ haben im Allgemeinen unterschiedliche $Q$-Werte am selben PV-Knoten. Die Matrix $\dotS$ ändert sich daher sowohl über die Iterationen szenarioabhängig als auch nach jeder Q-Korrektur.
%\end{enumerate}

%\subsection{Mathematische Darstellung des Problems}
%\label{subsec:tpf_pv_math}
%
%Die Iteration für den Spezialfall $\alpha_P = 1$ kann geschrieben werden als:
%\begin{equation}
%	\dotV_{(n+1)} = \ZB \cdot \underbrace{\left(\dotS^{*(n)} \had \dotV_{(n)}^{*\eleminv}\right)}_{\text{iterationsabhängig durch } Q^{(n)}} + \dot{W}
%	\label{eq:tpf_pv_iteration}
%\end{equation}
%
%wobei $\dotS^{*(n)}$ den konjugierten Leistungsvektor in Iteration $n$ bezeichnet, der sich für PV-Knoten ändert:
%\begin{equation}
%	\dotS^{*(n)}[k,:] = P_k^{\spec} - \jj Q_k^{(n)} \quad \text{für PV-Knoten } k
%	\label{eq:S_iterationsabh}
%\end{equation}
%
%Es ist wichtig festzuhalten, dass die Matrix $\ZB$ selbst weiterhin konstant bleibt -- sie hängt nur von der Netztopologie ab, nicht von den Lastdaten. Lediglich der \textit{Inhalt} der rechten Seite ändert sich. Dies ist ein wesentlicher Unterschied zum NR-Verfahren, bei dem sich die gesamte Jacobi-Matrix ändert, und bildet den Ansatzpunkt für die in \cref{ch:methoden_pv} vorgestellten Erweiterungen.

%\subsection{Zusammenfassende Gegenüberstellung}
%\label{subsec:tpf_gegenueber}
%
%\Cref{tab:tpf_pq_vs_pv} fasst die Auswirkung von PV-Knoten auf die TPF-Formulierung zusammen.
%
%\begin{table}[H]
%	\centering
%	\caption{Auswirkung von PV-Knoten auf den Tensor Power Flow.}
%	\label{tab:tpf_pq_vs_pv}
%	\begin{tabular}{@{}p{4.5cm}cc@{}}
%		\toprule
%		\textbf{Eigenschaft} & \textbf{Nur PQ-Knoten} & \textbf{Mit PV-Knoten} \\
%		\midrule
%		$\ZB$ konstant über Iterationen & \checkmark & \checkmark \\
%		$\vect{w}$ konstant über Iterationen & \checkmark & \checkmark \\
%		$\dotS$ konstant über Iterationen & \checkmark & $\times$ \\
%		$\dotS$ identisch für alle $\tau$ am gleichen Knoten & möglich & $\times$ \\
%		Tensor-Reshape anwendbar & \checkmark & Eingeschränkt \\
%		Banach-Konvergenzbeweis gültig & \checkmark & Nicht direkt \\
%		\bottomrule
%	\end{tabular}
%\end{table}

%Die Limitation des TPF auf PQ-Knoten ist in der Praxis für Verteilnetze oft akzeptabel, da dort die Mehrzahl der Einspeiser (PV-Anlagen mit Wechselrichter, Batteriespeicher) als PQ-Knoten mit festem Leistungsfaktor modelliert wird \cite{Salazar2024}. Mit zunehmender Integration von Synchrongeneratoren und spannungsregelnden Einheiten in Verteilnetzen gewinnt die korrekte Abbildung von PV-Knoten jedoch an Bedeutung. %Die Entwicklung geeigneter Erweiterungen des TPF zur Berücksichtigung von PV-Knoten ist Gegenstand von \cref{ch:methoden_pv}.





%
%\chapter{Der Tensor Power Flow}
%\label{sec:tpf}
%
%Aufbauend auf den in Abschnitt~\ref{sec:fpi} hergeleiteten Eigenschaften der Fixpunkt-Iteration wird in diesem Kapitel die Tensorisierung der FPI zum \emph{Tensor Power Flow} (TPF) nach Salazar Duque et al.~\cite{Salazar2024} entwickelt. Das zentrale Anliegen ist die Umformulierung der Iteration derart, dass $\tau$ Lastflüsse nicht sequenziell, sondern als eine einzige, hoch parallelisierbare Matrix-Operation gelöst werden. Der Abschnitt~\ref{subsec:tpf_motivation} motiviert diesen Übergang, die Abschnitte~\ref{subsec:tpf_dense} und \ref{subsec:tpf_sparse} entwickeln die Dense- bzw. Sparse-Formulierung, Abschnitt~\ref{subsec:tpf_komplexitaet} analysiert die Rechenkomplexität und Abschnitt~\ref{subsec:tpf_grenzen} zeigt die konzeptionelle Grenze des Standard-TPF auf, die die Notwendigkeit der in Kap.~\ref{sec:pv_methoden} entwickelten PV-Erweiterungen begründet.
%
%\section{Motivation und Konzept der Tensorisierung}
%\label{subsec:tpf_motivation}
%
%Die in Abschnitt~\ref{sec:fpi} hergeleitete FPI \eqref{eq:fpi_zbus} berechnet eine einzelne Spannungslösung $\mathbf{v}_d \in \mathbb{C}^{b \phi}$ für einen einzelnen Leistungsvektor $\mathbf{s}_d \in \mathbb{C}^{b\phi}$. In den in dieser Arbeit betrachteten multidimensionalen Anwendungen, wie Zeitreihensimulationen, probabilistischen Lastflüssen und Sensitivitätsstudien, sind jedoch typischerweise $\tau \gg 1$ Szenarien zu berechnen, was zu einer Datenstruktur der Form $\mathcal{S} \in \mathbb{C}^{b\phi \times \tau}$ führt. Die naive Anwendung der FPI besteht darin, für jede der $\tau$ Spalten von $\mathcal{S}$ die Iteration \eqref{eq:fpi_zbus} separat durchzuführen, was einer sequenziellen Schleife über $\tau$ entspricht. Diese Vorgehensweise verschenkt jedoch die zentrale strukturelle Eigenschaft der FPI: die Konstanz von $\mathbf{Z}_B$ über alle Szenarien.
%
%Das Konzept der Tensorisierung besteht darin, die $\tau$ einzelnen Matrix-Vektor-Multiplikationen $\mathbf{Z}_B \mathbf{x}_i$ mit $\mathbf{x}_i \in \mathbb{C}^{b\phi}$ ($i = 1,\ldots,\tau$) durch \emph{eine} Matrix-Matrix-Multiplikation $\mathbf{Z}_B \mathbf{X}$ mit $\mathbf{X} \in \mathbb{C}^{b\phi \times \tau}$ zu ersetzen. Der entscheidende Vorteil liegt darin, dass Matrix-Matrix-Multiplikationen im Vergleich zu Matrix-Vektor-Multiplikationen drastisch bessere Cache-Lokalität aufweist und somit die Rechenzeit reduziert~\cite{Salazar2024}.
%
%Für allgemeinere Studiendesigns mit mehreren unabhängigen Dimensionen (z.\,B.\ $p$ Experimente, $r$ Szenarien pro Experiment, $t$ Zeitschritte) lässt sich die Leistungsdatenstruktur als Tensor
%\begin{equation}
%	\label{eq:power_tensor}
%	\mathcal{S} \in \mathbb{C}^{\ldots \times p \times r \times t \times b\phi}
%\end{equation}
%darstellen. Alle unabhängigen Dimensionen können jedoch für die Berechnung zu einer einzigen effektiven Dimension $\tau = \ldots \cdot p \cdot r \cdot t$ zusammengefasst werden, ohne dass die Konvergenzeigenschaften der FPI berührt werden~\cite{Salazar2024}. Diese Reshape-Operation ist in modernen numerischen Bibliotheken (NumPy) speichereffizient möglich und wird im Folgenden als Ausgangspunkt für die Dense- und Sparse-Formulierung verwendet.
%
%\section{Dense-Formulierung}
%\label{subsec:tpf_dense}
%
%\subsubsection{Herleitung der Batched-Iterationsvorschrift}
%\label{subsubsec:tpf_dense_iteration}
%
%Ausgangspunkt ist die skalare FPI-Vorschrift \eqref{eq:fpi_zbus} für konstante Leistung:
%\begin{equation}
%	\label{eq:fpi_skalar_wiederholung}
%	\mathbf{v}_d^{(n+1)} = \mathbf{Z}_B \left( \mathbf{s}_d^{*} \oslash \mathbf{v}_d^{*\,(n)} \right) + \mathbf{w}
%\end{equation}
%mit der konstanten Bus-Impedanzmatrix $\mathbf{Z}_B = \mathbf{Y}_{dd}^{-1}$ und dem konstanten Vektor $\mathbf{w} = -\mathbf{Z}_B \mathbf{Y}_{ds} \mathbf{v}_s$. Die Iteration wird nun für $\tau$ Szenarien simultan formuliert, indem die Vektoren durch Matrizen ersetzt werden, in denen die $\tau$ Szenarien entlang der zweiten Dimension angeordnet sind:
%\begin{equation}
%	\dot{\mathbf{S}} = \left[ \mathbf{s}_d^{(1)}, \mathbf{s}_d^{(2)}, \ldots, \mathbf{s}_d^{(\tau)} \right] \in \mathbb{C}^{b\phi \times \tau}, \quad
%	\dot{\mathbf{V}} = \left[ \mathbf{v}_d^{(1)}, \mathbf{v}_d^{(2)}, \ldots, \mathbf{v}_d^{(\tau)} \right] \in \mathbb{C}^{b\phi \times \tau}
%\end{equation}
%Die Punktnotation kennzeichnet dabei die reshapte Tensorform. Die batched Iterationsvorschrift lautet somit
%\begin{equation}
%	\label{eq:tpf_dense}
%	\boxed{\;\dot{\mathbf{V}}^{(n+1)} = \mathbf{Z}_B \left( \dot{\mathbf{S}}^{*} \oslash \dot{\mathbf{V}}^{*\,(n)} \right) + \mathbf{W}\;}
%\end{equation}
%wobei $\mathbf{W} = \mathbf{w} \mathbf{1}_\tau^\top \in \mathbb{C}^{b\phi \times \tau}$ die spaltenweise Replikation des konstanten Slack-Beitrags über alle $\tau$ Szenarien darstellt. Alle Operationen in \eqref{eq:tpf_dense} sind elementweise bzw.\ als Matrix-Matrix-Multiplikation formuliert und lassen sich direkt auf BLAS-Routinen abbilden.
%
%Es sei ausdrücklich betont, dass die Reshape-Operation und die Batch-Auswertung die Konvergenzeigenschaften der FPI \emph{nicht} beeinflussen: Jede Spalte von $\dot{\mathbf{V}}$ folgt genau derselben skalaren Iteration \eqref{eq:fpi_skalar_wiederholung}, sodass die in Abschnitt~\ref{subsec:banach} bewiesene Konvergenzaussage für jede einzelne Spalte separat gilt~\cite{salazar2024}.
%
%\subsubsection{Algorithmische Umsetzung}
%\label{subsubsec:tpf_dense_algo}
%
%Algorithmus~\ref{alg:tpf_dense} zeigt die konkrete Implementierung der Dense-Formulierung. Die drei wesentlichen Kostenanteile sind (i) die einmalige Inversion von $\mathbf{Y}_{dd}$ zu Beginn, (ii) die Batched Hadamard-Division im Zähler und (iii) die zentrale Matrix-Matrix-Multiplikation $\mathbf{Z}_B \cdot (\cdot)$ pro Iteration.
%
%\begin{algorithm}[H]
%	\caption{Tensor Power Flow -- Dense-Formulierung}
%	\label{alg:tpf_dense}
%	\begin{algorithmic}[1]
%		\STATE \textbf{Eingabe:} $\dot{\mathbf{S}} \in \mathbb{C}^{b\phi \times \tau}$, $\mathbf{Y}_{dd}, \mathbf{Y}_{ds}$, $\mathbf{v}_s$, Toleranz $\varepsilon$, max.\ Iterationen $N_{\max}$
%		\STATE \textbf{Ausgabe:} $\dot{\mathbf{V}}^{(n)} \in \mathbb{C}^{b\phi \times \tau}$, Iterationszahl $n$
%		\STATE $\mathbf{Z}_B \leftarrow \mathbf{Y}_{dd}^{-1}$ \COMMENT{einmalige Vorabberechnung}
%		\STATE $\mathbf{W} \leftarrow -\mathbf{Z}_B \mathbf{Y}_{ds} \mathbf{v}_s \mathbf{1}_\tau^\top$ \COMMENT{Broadcasting}
%		\STATE $\dot{\mathbf{V}}^{(0)} \leftarrow \mathbf{1}_{b\phi \times \tau}$ \COMMENT{Flat-Start}
%		\STATE $n \leftarrow 0$, $\text{tol} \leftarrow \infty$
%		\WHILE{$\text{tol} \geq \varepsilon$ \textbf{ und } $n < N_{\max}$}
%		\STATE $\dot{\mathbf{V}}^{(n+1)} \leftarrow \mathbf{Z}_B \left( \dot{\mathbf{S}}^{*} \oslash \dot{\mathbf{V}}^{*\,(n)} \right) + \mathbf{W}$
%		\STATE $\text{tol} \leftarrow \max_{i,j} \left| \dot{V}^{(n+1)}_{ij} - \dot{V}^{(n)}_{ij} \right|$
%		\STATE $n \leftarrow n + 1$
%		\ENDWHILE
%		\STATE \textbf{return } $\dot{\mathbf{V}}^{(n)}, n$
%	\end{algorithmic}
%\end{algorithm}
%
%Bemerkenswert ist, dass die Iteration erst dann terminiert, wenn \emph{alle} $\tau$ Szenarien konvergiert sind. Dies erscheint zunächst als Nachteil, da einzelne Szenarien mit besonders schlechter Konditionierung ($\eta_i \to 1$) die Gesamtlaufzeit dominieren können. In der Praxis überwiegt jedoch der Vorteil der einheitlichen Batched-Operation deutlich, sofern die Szenarienvarianz moderat ist~\cite{salazar2024}. Alternativ lässt sich eine \emph{Masked Iteration} implementieren, bei der bereits konvergierte Spalten in nachfolgenden Iterationen nicht mehr aktualisiert werden -- dies wurde in der vorliegenden Arbeit als Implementierungsdetail übernommen, hat aber keinen Einfluss auf die theoretischen Konvergenzeigenschaften.
%
%\subsubsection{Speicheranforderung}
%\label{subsubsec:tpf_dense_speicher}
%
%Der wesentliche Nachteil der Dense-Formulierung liegt im Speicherbedarf der invertierten Systemmatrix. Während $\mathbf{Y}_{dd}$ typischerweise dünn besetzt ist mit $\mathcal{O}(b)$ Nicht-Null-Einträgen für radiale Netze, ist die inverse Matrix $\mathbf{Z}_B$ im Allgemeinen \emph{voll besetzt} mit $b^2 \phi^2$ komplexen Einträgen. Bei einfacher Genauigkeit (Complex64, 8 Byte pro Eintrag) ergibt sich der Speicherbedarf zu
%\begin{equation}
%	\label{eq:speicher_dense}
%	M_{\mathbf{Z}_B} = 8 \cdot (b\phi)^2 \text{ Byte}
%\end{equation}
%Für ein Netz mit $b\phi = 1000$ Bus-Phasen beträgt der Speicherbedarf $M_{\mathbf{Z}_B} \approx 8\,\mathrm{MB}$, was unkritisch ist; für $b\phi = 5000$ steigt er jedoch bereits auf $\approx 200\,\mathrm{MB}$ und für $b\phi = 10\,000$ auf $\approx 800\,\mathrm{MB}$. Diese quadratische Skalierung motiviert die im nächsten Abschnitt entwickelte Sparse-Formulierung.
%
%\section{Sparse-Formulierung}
%\label{subsec:tpf_sparse}
%
%\subsubsection{Vermeidung der expliziten Inversion}
%\label{subsubsec:tpf_sparse_motivation}
%
%Die Grundidee der Sparse-Formulierung besteht darin, die explizite Berechnung von $\mathbf{Z}_B = \mathbf{Y}_{dd}^{-1}$ zu vermeiden und stattdessen in jedem Iterationsschritt ein lineares Gleichungssystem der Form $\mathbf{Y}_{dd} \mathbf{v}_d^{(n+1)} = \mathbf{b}^{(n)}$ zu lösen. Da $\mathbf{Y}_{dd}$ konstant ist und in typischen Verteilnetzen eine Sparsity von $\mathcal{O}(b)$ Nicht-Null-Einträgen aufweist, kann eine einmalige LU-Zerlegung
%\begin{equation}
%	\label{eq:lu_zerlegung}
%	\mathbf{Y}_{dd} = \mathbf{P} \mathbf{L} \mathbf{U} \mathbf{Q}
%\end{equation}
%mit Permutationsmatrizen $\mathbf{P}, \mathbf{Q}$ und dünn besetzten Dreiecksmatrizen $\mathbf{L}, \mathbf{U}$ vor Beginn der Iteration berechnet werden. Anschließend erfordert jede Iteration nur noch eine Vorwärts- und Rückwärtssubstitution, deren Aufwand für typische radiale Netztopologien $\mathcal{O}(n_{nz})$ mit $n_{nz}$ als Anzahl der Nicht-Null-Einträge in $\mathbf{L}\mathbf{U}$ beträgt~\cite{davis2006}.
%
%Salazar Duque et al.~\cite{salazar2024} formulieren die Sparse-Variante durch Umstellung der FPI \eqref{eq:fpi_linsystem}. Für den Fall konstanter Leistung ergibt sich unter Verwendung der Definitionen aus \eqref{eq:fpi_konstant}:
%\begin{equation}
%	\label{eq:tpf_sparse_vektor}
%	\dot{\mathbf{M}} \dot{\mathbf{V}}^{(n+1)} = \dot{\mathbf{V}}^{*\,(n)\,\circ(-1)} + \dot{\mathbf{H}}
%\end{equation}
%mit den szenarienspezifischen Sparse-Blöcken
%\begin{equation}
%	\label{eq:m_h_bloecke}
%	\mathbf{M}^{(i)} = -\operatorname{diag}\left( \mathbf{s}_d^{*\,(i)\,\circ(-1)} \right) \mathbf{Y}_{dd}, \quad
%	\mathbf{h}^{(i)} = -\operatorname{diag}\left( \mathbf{s}_d^{*\,(i)\,\circ(-1)} \right) \mathbf{Y}_{ds} \mathbf{v}_s
%\end{equation}
%für $i = 1, \ldots, \tau$. Die Diagonalskalierung mit $\mathbf{s}_d^{*\,\circ(-1)}$ dient dazu, die Nichtlinearität aus dem Zähler in eine konstante Systemmatrix pro Szenario zu überführen.
%
%\subsubsection{Blockdiagonale Reshape-Struktur}
%\label{subsubsec:tpf_sparse_struktur}
%
%Um die $\tau$ separaten linearen Systeme in \emph{ein} großes Sparse-System zu überführen, werden die $\mathbf{M}^{(i)}$ als Blöcke einer blockdiagonalen Matrix arrangiert:
%\begin{equation}
%	\label{eq:blockdiagonal}
%	\dot{\mathbf{M}} = 
%	\begin{bmatrix}
%		\mathbf{M}^{(1)} & & & \\
%		& \mathbf{M}^{(2)} & & \\
%		& & \ddots & \\
%		& & & \mathbf{M}^{(\tau)}
%	\end{bmatrix}
%	\in \mathbb{C}^{(b\phi \cdot \tau) \times (b\phi \cdot \tau)}
%\end{equation}
%und die zugehörigen $\mathbf{h}^{(i)}$ sowie Spannungs- und Leistungsvektoren vertikal konkateniert:
%\begin{equation}
%	\dot{\mathbf{H}} = \begin{bmatrix} \mathbf{h}^{(1)} \\ \vdots \\ \mathbf{h}^{(\tau)} \end{bmatrix}, \quad
%	\dot{\mathbf{V}} = \begin{bmatrix} \mathbf{v}_d^{(1)} \\ \vdots \\ \mathbf{v}_d^{(\tau)} \end{bmatrix} \in \mathbb{C}^{b\phi \cdot \tau}
%\end{equation}
%Das resultierende Sparse-System hat die Gestalt eines linearen Gleichungssystems $\mathbf{A} \mathbf{x} = \mathbf{b}$, welches iterativ gelöst wird:
%\begin{equation}
%	\label{eq:tpf_sparse_system}
%	\boxed{\;\dot{\mathbf{M}} \dot{\mathbf{V}}^{(n+1)} = \dot{\mathbf{V}}^{*\,(n)\,\circ(-1)} + \dot{\mathbf{H}}\;}
%\end{equation}
%
%\paragraph{Bemerkung zur Sparsity-Struktur.}
%Da die einzelnen Blöcke $\mathbf{M}^{(i)}$ dieselbe Nichtnullstruktur wie $\mathbf{Y}_{dd}$ besitzen (die Diagonalskalierung ändert die Sparsity nicht), erbt $\dot{\mathbf{M}}$ die Blockdiagonalstruktur mit einer Gesamtzahl von $\tau \cdot n_{nz}(\mathbf{Y}_{dd})$ Nicht-Null-Einträgen. Dies wächst nur linear mit $\tau$, im Gegensatz zum quadratisch wachsenden Dense-Speicher \eqref{eq:speicher_dense}. Zudem lässt sich die symbolische Faktorisierung eines einzelnen $\mathbf{M}^{(i)}$-Blocks aufgrund der identischen Struktur vorab berechnen und auf alle $\tau$ Szenarien anwenden~\cite{davis2006}.
%
%\subsubsection{Algorithmische Umsetzung}
%\label{subsubsec:tpf_sparse_algo}
%
%Algorithmus~\ref{alg:tpf_sparse} zeigt die Sparse-Implementierung. Der zentrale Punkt ist, dass die symbolische Analyse (Fill-In-Reduktion durch AMD-, COLAMD- oder METIS-Ordering) und die numerische LU-Faktorisierung von $\dot{\mathbf{M}}$ nur \emph{einmal} vor Beginn der Iterationsschleife durchgeführt werden. Innerhalb der Schleife reduziert sich der Aufwand auf eine Vorwärts-/Rückwärtssubstitution, deren Kosten für radiale Netze $\mathcal{O}(b\phi \cdot \tau)$ beträgt~\cite{davis2006}.
%
%\begin{algorithm}[H]
%	\caption{Tensor Power Flow -- Sparse-Formulierung}
%	\label{alg:tpf_sparse}
%	\begin{algorithmic}[1]
%		\STATE \textbf{Eingabe:} $\dot{\mathbf{S}} \in \mathbb{C}^{b\phi \times \tau}$, sparse $\mathbf{Y}_{dd}, \mathbf{Y}_{ds}$, $\mathbf{v}_s$, Toleranz $\varepsilon$, max.\ Iterationen $N_{\max}$
%		\STATE \textbf{Ausgabe:} $\dot{\mathbf{V}}^{(n)}$, Iterationszahl $n$
%		\STATE $\dot{\mathbf{M}} \leftarrow \text{blockdiag}\left(\left\{ -\operatorname{diag}(\mathbf{s}_d^{*\,(i)\,\circ(-1)}) \mathbf{Y}_{dd} \right\}_{i=1}^{\tau}\right)$
%		\STATE $\dot{\mathbf{H}} \leftarrow \text{vstack}\left(\left\{ -\operatorname{diag}(\mathbf{s}_d^{*\,(i)\,\circ(-1)}) \mathbf{Y}_{ds} \mathbf{v}_s \right\}_{i=1}^{\tau}\right)$
%		\STATE $(\mathbf{P}, \mathbf{L}, \mathbf{U}, \mathbf{Q}) \leftarrow \text{LU-factorize}(\dot{\mathbf{M}})$ \COMMENT{einmalige Faktorisierung}
%		\STATE $\dot{\mathbf{V}}^{(0)} \leftarrow \mathbf{1}_{b\phi \cdot \tau}$
%		\STATE $n \leftarrow 0$, $\text{tol} \leftarrow \infty$
%		\WHILE{$\text{tol} \geq \varepsilon$ \textbf{ und } $n < N_{\max}$}
%		\STATE $\mathbf{b}^{(n)} \leftarrow \dot{\mathbf{V}}^{*\,(n)\,\circ(-1)} + \dot{\mathbf{H}}$
%		\STATE $\dot{\mathbf{V}}^{(n+1)} \leftarrow \text{solve\_LU}(\mathbf{P}, \mathbf{L}, \mathbf{U}, \mathbf{Q}, \mathbf{b}^{(n)})$
%		\STATE $\text{tol} \leftarrow \max_i \left| \dot{V}^{(n+1)}_i - \dot{V}^{(n)}_i \right|$
%		\STATE $n \leftarrow n + 1$
%		\ENDWHILE
%		\STATE $\dot{\mathbf{V}}^{(n)} \leftarrow \text{reshape}\left(\dot{\mathbf{V}}^{(n)}, (b\phi, \tau)\right)$
%		\STATE \textbf{return } $\dot{\mathbf{V}}^{(n)}, n$
%	\end{algorithmic}
%\end{algorithm}
%
%\paragraph{Implementierungsdetail.}
%In der Praxis kann die blockdiagonale Struktur von $\dot{\mathbf{M}}$ ausgenutzt werden, um die Sparse-Faktorisierung parallel über die $\tau$ Blöcke zu verteilen: Jeder Block $\mathbf{M}^{(i)}$ wird unabhängig faktorisiert, was für $\tau \gg 1$ zu einer weiteren Beschleunigung führt. In der vorliegenden Implementierung wird dieser Ansatz mittels \texttt{scipy.sparse.linalg.splu} für Einzelblock-Faktorisierungen und einer batched Vorwärts-/Rückwärtssubstitution realisiert.
%
%\subsection{Rechenkomplexität und Vergleich Dense vs.\ Sparse}
%\label{subsec:tpf_komplexitaet}
%
%Die Rechenkomplexität der beiden TPF-Varianten pro Iteration setzt sich aus zwei Beiträgen zusammen: (i) den elementweisen Operationen (Hadamard-Division, Update) mit Aufwand $\mathcal{O}(b\phi \cdot \tau)$, welche in beiden Varianten identisch sind, und (ii) der zentralen linearen Operation, welche sich substanziell unterscheidet.
%
%\paragraph{Dense-Formulierung.}
%Die Matrix-Matrix-Multiplikation $\mathbf{Z}_B \cdot \mathbf{X}$ mit $\mathbf{Z}_B \in \mathbb{C}^{b\phi \times b\phi}$ und $\mathbf{X} \in \mathbb{C}^{b\phi \times \tau}$ hat den Aufwand
%\begin{equation}
%	\label{eq:komplex_dense}
%	C_{\text{dense}}^{\text{iter}} = \mathcal{O}\left( (b\phi)^2 \cdot \tau \right)
%\end{equation}
%Vorabkosten für die Inversion: $\mathcal{O}\left((b\phi)^3\right)$. Der Speicheraufwand skaliert quadratisch mit $b\phi$.
%
%\paragraph{Sparse-Formulierung.}
%Für radiale Netze mit $n_{nz}(\mathbf{L}\mathbf{U}) = \mathcal{O}(b\phi)$ (dank Fill-In-reduzierendem Ordering) beträgt der Aufwand der Vorwärts-/Rückwärtssubstitution pro Iteration
%\begin{equation}
%	\label{eq:komplex_sparse}
%	C_{\text{sparse}}^{\text{iter}} = \mathcal{O}\left( n_{nz}(\mathbf{L}\mathbf{U}) \cdot \tau \right) = \mathcal{O}\left( b\phi \cdot \tau \right)
%\end{equation}
%Vorabkosten für die LU-Faktorisierung: $\mathcal{O}\left((b\phi)^{1{,}4}\right)$ für typische Netztopologien~\cite{davis2006}. Der Speicheraufwand skaliert linear mit $b\phi$.
%
%Tabelle~\ref{tab:tpf_vergleich} fasst die charakteristischen Eigenschaften der beiden Formulierungen zusammen.
%
%\begin{table}[H]
%	\centering
%	\caption{Vergleich der Dense- und Sparse-Formulierung des TPF}
%	\label{tab:tpf_vergleich}
%	\begin{tabular}{lll}
%		\toprule
%		Eigenschaft & Dense-Formulierung & Sparse-Formulierung \\
%		\midrule
%		Systemmatrix & $\mathbf{Z}_B = \mathbf{Y}_{dd}^{-1}$ (voll) & $\mathbf{Y}_{dd}$ (dünn) \\
%		Vorabaufwand & $\mathcal{O}((b\phi)^3)$ Inversion & $\mathcal{O}((b\phi)^{1{,}4})$ LU \\
%		Aufwand pro Iteration & $\mathcal{O}((b\phi)^2 \cdot \tau)$ & $\mathcal{O}(b\phi \cdot \tau)$ \\
%		Speicheraufwand & $\mathcal{O}((b\phi)^2)$ & $\mathcal{O}(b\phi \cdot \tau)$ \\
%		Skalierung mit $b\phi$ & quadratisch & linear \\
%		BLAS-Nutzung & Level 3 (\texttt{GEMM}) & Level 2 (Substitution) \\
%		GPU-Eignung & ausgezeichnet & moderat \\
%		Empfohlener Bereich & $b\phi \lesssim 2000$ & $b\phi \gtrsim 2000$ \\
%		\bottomrule
%	\end{tabular}
%\end{table}
%
%\paragraph{Empirische Break-Even-Analyse.}
%Salazar Duque et al.~\cite{salazar2024} untersuchen empirisch den Übergang zwischen Dense- und Sparse-Präferenz. Für $\tau = 500$ Lastflüsse liegt der Break-Even bei ca.\ $b\phi \approx 900$ Bus-Phasen: Für kleinere Netze ist die Dense-Formulierung dank optimierter BLAS-\texttt{GEMM}-Routinen überlegen, für größere Netze überwiegt der Sparsity-Vorteil. Diese Beobachtung wird in Kap.~\ref{sec:ergebnisse} für die im Rahmen dieser Arbeit implementierten Varianten (mit PV-Erweiterung) verifiziert.
%
%\subsection{GPU-Beschleunigung (optional)}
%\label{subsec:tpf_gpu}
%
%Die Dense-Formulierung eignet sich aufgrund ihrer Struktur (Batched-\texttt{GEMM}) hervorragend für die Ausführung auf Grafikprozessoren. Moderne GPUs bieten typischerweise eine Peak-Performance im Bereich von $10\text{--}20\,\mathrm{TFLOPS}$ für einfach-präzise komplexe Arithmetik, was einer Beschleunigung um den Faktor $10\text{--}50$ gegenüber CPU-basierten Implementierungen entspricht. Salazar Duque et al.~\cite{salazar2024} zeigen für eine Jahressimulation im Minutenraster ($\tau = 525\,600$) auf einem 100-Bus-Netz eine Reduktion der Rechenzeit von $22\,\mathrm{s}$ (CPU Dense) auf $78\,\mathrm{s}$ (GPU Dense), wobei der GPU-Overhead durch Daten-Transfer zwischen Host und Device für kleine $\tau$ die CPU-Version zunächst unterlegen macht. Für $\tau \gg 10^4$ und $b\phi \gtrsim 10^3$ dominiert jedoch die GPU-Version.
%
%In der vorliegenden Arbeit ist die GPU-Unterstützung als optionale Erweiterung mittels CuPy vorgesehen, jedoch nicht zentraler Untersuchungsgegenstand. Der Fokus liegt auf der methodischen Erweiterung des CPU-basierten TPF um PV-Knoten.
%
%\subsection{Grenzen des Standard-TPF und Ausblick auf PV-Integration}
%\label{subsec:tpf_grenzen}
%
%Die in diesem Kapitel dargestellte TPF-Formulierung, sowohl in ihrer Dense- als auch in ihrer Sparse-Variante, beruht fundamental auf zwei Annahmen:
%
%\begin{enumerate}
%	\item \textbf{Konstante Leistungsvorgabe an allen Lastknoten.} Der Leistungsvektor $\mathbf{s}_d = \mathbf{P}_d + j\mathbf{Q}_d$ ist vollständig vorgegeben und ändert sich weder innerhalb einer Iteration noch zwischen Iterationen.
%	\item \textbf{Konstanz der Systemmatrix.} Die Matrix $\mathbf{Z}_B$ (Dense) bzw.\ $\mathbf{Y}_{dd}$ (Sparse) hängt ausschließlich von der Netztopologie und den Leitungsparametern ab, nicht vom Iterationszustand.
%\end{enumerate}
%
%Beide Annahmen sind für Netze mit ausschließlich PQ-Knoten und Slack-Bus erfüllt und ermöglichen die massive Parallelisierung, die den Kernvorteil des TPF ausmacht. Sobald jedoch \emph{Spannungsgeregelte Knoten (PV-Knoten)} im Netz vorhanden sind -- typischerweise durch Generatoren, Synchronkondensatoren oder wechselrichterbasierte Erzeugungsanlagen mit Spannungsregelung --, werden beide Annahmen verletzt:
%
%\begin{itemize}
%	\item An PV-Knoten $k \in \Omega_{PV}$ ist der Spannungsbetrag $|V_k| = |V_k^{\text{spec}}|$ vorgegeben, während die Blindleistung $Q_k$ eine unbekannte Zustandsgröße darstellt. Der entsprechende Eintrag im Leistungsvektor $\mathbf{s}_d$ ist somit nicht mehr vollständig bekannt, sondern muss iterativ mitbestimmt werden.
%	\item Die Kopplung zwischen $Q_k$ und dem Spannungszustand über die nichtlineare Beziehung $|V_k(\mathbf{Q})| = |V_k^{\text{spec}}|$ zerstört die Konstanz der effektiven Systemmatrix, da jede Q-Aktualisierung eine Änderung des Nichtlinearitätsterms $\mathbf{s}^{*} \oslash \mathbf{v}^{*}$ zur Folge hat, welche wiederum eine Anpassung der Iteration erfordert.
%\end{itemize}
%
%Diese strukturelle Inkompatibilität von PV-Knoten mit der ursprünglichen TPF-Formulierung ist der methodische Kern der vorliegenden Arbeit. Salazar Duque et al.~\cite{Salazar2024} schließen ihre Formulierung explizit auf PQ-Knoten ein und weisen darauf hin, dass die Behandlung von PV-Knoten im Rahmen ihrer FPI-Systematik einer eigenständigen Untersuchung bedürfe. Auch in der breiteren Literatur zur strombasierten Lastflussberechnung wurde diese Problematik früh erkannt: Tinney~\cite{tinney1991} bemerkte, dass strombasierte Verfahren mit konstanter Admittanzmatrix keine allgemeine PV-Behandlung besitzen und die bekannten Ansätze zur Einbindung von PV-Knoten die Konvergenzgeschwindigkeit zunichtemachen. da Costa et al.~\cite{daCosta1999} lieferten für das strombasierte NR-Verfahren eine algorithmisch elegante Lösung durch die Einführung einer zusätzlichen Variablen $\Delta Q$ pro PV-Knoten (vgl.\ Kap.~\ref{subsec:nr_pv}); im FPI-Kontext ist eine solche Erweiterung jedoch nicht direkt übertragbar, da die FPI keine Jacobi-Matrix aufbaut, in die zusätzliche Zeilen und Spalten integriert werden könnten.
%
%Das folgende Kapitel~\ref{sec:pv_methoden} entwickelt daher zwei alternative Ansätze, die die Parallelisierbarkeit des TPF weitestgehend erhalten:
%
%\begin{itemize}
%	\item \textbf{Methode A (Äußere Q-Schleife):} Der Standard-TPF wird als innere Schleife mit fest vorgegebenem $\mathbf{Q}_{PV}$ ausgeführt; eine äußere Schleife aktualisiert die Blindleistungen basierend auf der Spannungsabweichung und einer diagonalen Approximation der Sensitivitätsmatrix $\partial |V|^2 / \partial Q$. Die Innere-Schleifen-Struktur bleibt vollständig tensorisierbar.
%	\item \textbf{Methode B (Interleaved Update):} Nach jedem FPI-Schritt wird die Spannung an PV-Knoten auf den Sollbetrag projiziert und die Blindleistung direkt aus der Leistungsbilanz berechnet. Sämtliche Operationen bleiben elementweise und mit der Tensorstruktur kompatibel.
%\end{itemize}
%
%Beide Methoden werden in Dense- und Sparse-Formulierung implementiert und gegen die NR-Referenzlösung von \texttt{pandapower} validiert. Der zentrale Anspruch besteht darin, den in Tabelle~\ref{tab:nr_vergleich} dokumentierten Geschwindigkeitsvorteil des TPF auch für Netze mit signifikantem PV-Anteil weitgehend zu erhalten.
%
%\bigskip
%\noindent
%Zusammenfassend wurde in diesem Kapitel die Tensorisierung der FPI zum Tensor Power Flow in Dense- und Sparse-Formulierung entwickelt. Beide Varianten erhalten die Konvergenzeigenschaften der zugrundeliegenden FPI vollständig und ermöglichen die simultane Lösung von $\tau$ Lastflüssen mittels BLAS-basierter Batched-Operationen. Die Dense-Variante ist für kleine bis mittlere Netze ($b\phi \lesssim 2000$) überlegen, die Sparse-Variante skaliert dank ihrer linearen Speicheranforderung deutlich besser für große Netze. Die strukturelle Beschränkung auf reine PQ-Netze motiviert die im folgenden Kapitel entwickelten PV-Erweiterungen.